%%%%%%%%%%%%%%%%%%%%%%%%%%%%%%%%%%%%%%%%%%%%%%%%%%%%%%%%%%%%%%%%%%%%%%%%%%%%%%
%  Partie « SWOT et enjeux » (6 diapositives) -- à importer avec %%%%%%%%%%%%%%%%%%%%%%%%%%%%%%%%%%%%%%%%%%%%%%%%%%%%%%%%%%%%%%%%%%%%%%%%%%%%%%
%  Partie « SWOT et enjeux » (6 diapositives) -- à importer avec %%%%%%%%%%%%%%%%%%%%%%%%%%%%%%%%%%%%%%%%%%%%%%%%%%%%%%%%%%%%%%%%%%%%%%%%%%%%%%
%  Partie « SWOT et enjeux » (6 diapositives) -- à importer avec %%%%%%%%%%%%%%%%%%%%%%%%%%%%%%%%%%%%%%%%%%%%%%%%%%%%%%%%%%%%%%%%%%%%%%%%%%%%%%
%  Partie « SWOT et enjeux » (6 diapositives) -- à importer avec \include{chap_swot}
%  Fichier de contenu : ni préambule ni \begin{document}, commence par \section.
%
%  Remplace l'ancien SWOT en PDF (\includepdf supprimé) : le diagnostic est maintenant écrit
%  en LaTeX/TikZ, avec la charte du reste de la présentation, et complété par un SWOT croisé.
%  Ce qui a changé par rapport à la version PDF :
%   - la décomposition des 69 € n'est plus répétée ici (elle est dans la partie Finances) ;
%   - chaque point du diagnostic renvoie à un « frein » F1 à F7 de l'étude documentaire ;
%   - les marges des offres Solo / Duo / Groupe sont recalculées (marge par commande ET par kit,
%     avec le coût de vente de 4 € compté une seule fois par commande) ;
%   - le SWOT croisé (stratégies SO, WO, ST, WT) relie le diagnostic aux trois enjeux.
%  Hypothèses chiffrées : celles de la partie Finances (kit 69 € TTC, matériel 30,30 €, vente 4 €
%  par commande, 1 kit offert pour 10, charges fixes de l'année 3 = 430 k€).
%  Les offres Solo (79 €), Duo (139 €, 2 kits) et Groupe (249 €, 4 kits) sont celles du SWOT d'origine ;
%  le nombre de kits du Groupe (4) est une hypothèse.
%%%%%%%%%%%%%%%%%%%%%%%%%%%%%%%%%%%%%%%%%%%%%%%%%%%%%%%%%%%%%%%%%%%%%%%%%%%%%%
\section{SWOT et enjeux}

% Entrée de SWOT : \swit[couleur du code]{code}{titre}{texte}  (titre et texte sur le même paragraphe)
\newcommand{\swit}[4][insarouge]{%
  \begin{tcolorbox}[colback=insabeige,colframe=insabeige,boxrule=0pt,arc=2pt,left=4pt,right=4pt,
                    top=1.5pt,bottom=1.5pt,before skip=0pt,after skip=2pt]
    \scriptsize\raggedright{\sffamily\bfseries\color{#1}#2}\hspace{0.3em}{\sffamily\bfseries\footnotesize #3}\hspace{0.3em}#4\par
  \end{tcolorbox}}
% En-tête de colonne : \swhead{Forces}{couleur}
\newcommand{\swhead}[2]{%
  {\sffamily\bfseries\footnotesize\color{#2}\MakeUppercase{#1}}\\[-3pt]{\color{#2}\rule{\linewidth}{0.6pt}}\\[-1pt]}
% renvoi vers un enjeu ou un frein
\newcommand{\swref}[1]{\textcolor{insarouge}{\textbf{→\,#1}}}

%---------------------------------------------------------------------------
\begin{frame}{Diagnostic interne~: forces et faiblesses}
  \begin{columns}[T]
    \begin{column}{0.485\textwidth}
      \swhead{Forces}{insarouge}
      \swit{S1}{Proposition de valeur claire}{Une messagerie texte sans réseau mobile, sans compte ni abonnement.}
      \swit{S2}{Coût d'entrée faible}{69\,€ une fois, contre ≈ 830\,€ sur trois ans pour une balise satellite.}
      \swit{S3}{Écosystème libre existant}{Meshtastic, Heltec V3, applications~: l'effort porte sur l'usage.}
      \swit{S4}{Simplicité comme différence}{Réglages Europe, clé privée, QR code, guides en français.}
      \swit{S5}{Un socle, plusieurs usages}{Outdoor, campagne, ONG, crise~: un kit, quatre terrains.}
    \end{column}
    \begin{column}{0.485\textwidth}
      \swhead{Faiblesses}{insagris}
      \swit[insagris]{W1}{Faible barrière technologique}{Un bricoleur refait un kit proche pour 18 à 30\,€. \swref{F5}}
      \swit[insagris]{W2}{Effet de réseau indispensable}{Un kit isolé ne sert à rien~: la valeur vient des voisins. \swref{enjeu 1}}
      \swit[insagris]{W3}{Marge de contribution fine}{20,17\,€ par kit, mais 0,62\,€ de résultat par kit en année 3. \swref{enjeu 2}}
      \swit[insagris]{W4}{Financement avant rentabilité}{≈ 450\,k€ à financer, point mort ≈ 21\,300 kits par an. \swref{F7}}
      \swit[insagris]{W5}{Positionnement trop large}{ONG, randonneurs, makers, public~: attentes différentes. \swref{enjeu 3}}
    \end{column}
  \end{columns}
  \vspace{-0.2cm}
  \begin{pcompact}\centering\footnotesize
    Lecture clé~: la technologie est accessible~; la valeur viendra de l'usage, de la marque et de l'écosystème.
  \end{pcompact}
  \srcnote{Synthèse des parties Concurrence, Finances et Étude documentaire. F2 à F7 = freins de l'étude documentaire.}
\end{frame}

%---------------------------------------------------------------------------
\begin{frame}{Diagnostic externe~: opportunités et menaces}
  \begin{columns}[T]
    \begin{column}{0.485\textwidth}
      \swhead{Opportunités}{insarouge}
      \swit{O1}{Un marché outdoor important}{27 millions de pratiquants, dont 240\,000 licenciés FFRandonnée.}
      \swit{O2}{Besoin croissant de résilience}{Zones blanches, crises~: l'État invite à préparer un kit de 72\,h.}
      \swit{O3}{Réseaux locaux à construire}{Refuges, clubs, communes, ADRASEC~: points d'ancrage du maillage.}
      \swit{O4}{Précommandes adaptées}{Le financement participatif teste la demande et paie le premier stock.}
      \swit{O5}{Made in France, souveraineté}{Relocalisation par étapes~: un argument pour les collectivités.}
    \end{column}
    \begin{column}{0.485\textwidth}
      \swhead{Menaces}{insagris}
      \swit[insagris]{T1}{Fait-maison, kits concurrents}{Cartes dès 18\,€, kits prêts à l'emploi 59 à 153\,€. \swref{C}}
      \swit[insagris]{T2}{Substituts puissants}{SOS satellite des smartphones, balises, SMS satellite. \swref{F2}}
      \swit[insagris]{T3}{Contraintes réglementaires}{Marquage CE, cybersécurité, responsabilité produit. \swref{F3, F4}}
      \swit[insagris]{T4}{Dépendance amont}{Un seul fabricant de carte~; cap et marque de Meshtastic. \swref{F5, F6}}
      \swit[insagris]{T5}{Promesse mal comprise}{Ni téléphone, ni balise de détresse~: cadrer la promesse. \swref{F3}}
    \end{column}
  \end{columns}
  \vspace{-0.22cm}
  \begin{pcompact}\centering\footnotesize
    Opportunité réelle, mais fenêtre concurrentielle courte~: il faut bâtir vite une communauté et une marque.
  \end{pcompact}
  \srcnote{Sources~: \cite{ffrando,ffrando_licencies,guide_tous_resp}. F2 à F6 = freins de l'étude documentaire~; C = chapitre Concurrence.}
\end{frame}

%---------------------------------------------------------------------------
\newenvironment{pucesx}
  {\begingroup\setbeamerfont{itemize/enumerate body}{size=\scriptsize}\begin{itemize}\setlength{\itemsep}{0pt}\setlength{\topsep}{1pt}}
  {\end{itemize}\endgroup}
\begin{frame}{SWOT croisé~: quelles stratégies en déduire~?}
  \footnotesize\renewcommand{\arraystretch}{1.0}\setlength{\tabcolsep}{4pt}
  \hyphenpenalty=10000 \exhyphenpenalty=10000
  \noindent\begin{tabularx}{\textwidth}{@{}>{\raggedright\arraybackslash}p{1.65cm}>{\raggedright\arraybackslash}X>{\raggedright\arraybackslash}X@{}}
    & \cellcolor{insarose}\textcolor{insarouge}{\bfseries Opportunités (O)} & \cellcolor{insarose}\textcolor{insarouge}{\bfseries Menaces (T)} \\[2pt]
    \cellcolor{insarose}\textcolor{insarouge}{\bfseries Forces (S)} &
      \cellcolor{insabeige}{\sffamily\bfseries\color{insarouge}Attaquer}
      \begin{pucesx}
        \item Micro-réseaux prêts à l'emploi (Duo, Groupe) pour clubs, refuges, groupes de randonnée.
        \item Précommandes~: tester la demande et financer le premier stock.
        \item Pilotes avec des ADRASEC et des communes~: preuves d'usage, premiers nœuds.
      \end{pucesx} &
      \cellcolor{insabeige}{\sffamily\bfseries\color{insarouge}Défendre}
      \begin{pucesx}
        \item Vendre le prêt-à-l'emploi~: réglages, clé privée, guide, garantie, marquage CE.
        \item Rester un complément du smartphone et de la balise, sans promesse de secours.
        \item Contribuer au code libre, rester « compatible Meshtastic ».
      \end{pucesx} \\[3pt]
    \cellcolor{insarose}\textcolor{insarouge}{\bfseries Faiblesses (W)} &
      \cellcolor{insabeige}{\sffamily\bfseries\color{insarouge}Corriger}
      \begin{pucesx}
        \item Densifier autour de points d'ancrage (refuge, club, mairie) contre l'effet de réseau.
        \item Tester Solo, Duo et Groupe avant la série pour relever la marge.
        \item Bâtir le dossier de financement sur les précommandes.
      \end{pucesx} &
      \cellcolor{insabeige}{\sffamily\bfseries\color{insarouge}Éviter}
      \begin{pucesx}
        \item Notice claire, test à 100\,\%, assurance responsabilité produit.
        \item Double source de cartes (V3 et V4), stock tampon, vente d'abord en France.
        \item Fixer des seuils d'arrêt avant le moule et la série (étude quantitative).
      \end{pucesx} \\
  \end{tabularx}
  \srcnote{Croisement des deux diapositives précédentes (méthode TOWS)~: S, W = forces et faiblesses~; O, T = opportunités et menaces.}
\end{frame}

%---------------------------------------------------------------------------
% Liaisons possibles entre n kits : n(n-1)/2 (2 kits = 1, 4 = 6, 6 = 15, 10 = 45)
\begin{frame}{Enjeu 1~: résoudre l'effet de réseau}
  \begin{columns}[T]
    \begin{column}{0.46\textwidth}
      \centering
      \begin{tikzpicture}[font=\footnotesize]
        \node[pcli,text width=2.3cm,minimum height=0.7cm] (a) at (0,1.05) {Peu d'utilisateurs};
        \node[pgris,text width=1.6cm,minimum height=0.7cm] (b) at (3.55,1.05) {Peu de relais};
        \node[pgris,text width=1.6cm,minimum height=0.7cm] (c) at (1.78,-0.3) {Valeur limitée};
        \draw[fluxr] (a) -- (b);
        \draw[fluxr] (b) -- (c);
        \draw[fluxr] (c) -- (a);
        \node[font=\scriptsize\bfseries,text=insarouge,align=center] at (1.78,0.5) {cercle\\vicieux};
      \end{tikzpicture}

      \vspace{0.05cm}
      {\footnotesize\textbf{Un groupe vaut plus que la somme de ses kits}~: $n$ kits offrent $n(n{-}1)/2$ liaisons.\par}
      \begin{tikzpicture}[font=\scriptsize]
        \tikzset{noeud/.style={circle,fill=insarouge,inner sep=1.3pt}}
        % 2 kits
        \begin{scope}[shift={(0.0,0)}]
          \node[noeud] (u0) at (-0.28,0) {};\node[noeud] (u1) at (0.28,0) {};
          \draw[insagris,thick] (u0) -- (u1);
          \node[align=center,anchor=north] at (0,-0.55) {2 kits\\\textbf{1 liaison}};
        \end{scope}
        % 4 kits
        \begin{scope}[shift={(1.85,0)}]
          \foreach \i in {0,...,3} {\node[noeud] (v\i) at ({45+90*\i}:0.4) {};}
          \foreach \i/\j in {0/1,0/2,0/3,1/2,1/3,2/3} {\draw[insagris,thick] (v\i) -- (v\j);}
          \node[align=center,anchor=north] at (0,-0.55) {4 kits\\\textbf{6 liaisons}};
        \end{scope}
        % 6 kits
        \begin{scope}[shift={(3.7,0)}]
          \foreach \i in {0,...,5} {\node[noeud] (w\i) at ({90+60*\i}:0.42) {};}
          \foreach \i/\j in {0/1,0/2,0/3,0/4,0/5,1/2,1/3,1/4,1/5,2/3,2/4,2/5,3/4,3/5,4/5} {\draw[insagris,thin] (w\i) -- (w\j);}
          \node[align=center,anchor=north] at (0,-0.55) {6 kits\\\textbf{15 liaisons}};
        \end{scope}
        \node[align=center,font=\scriptsize] at (5.15,0) {10 kits\\\textbf{45}};
      \end{tikzpicture}
    \end{column}
    \begin{column}{0.51\textwidth}
      \begin{columns}[T]
        \begin{column}{0.47\linewidth}
          \begin{carte}[height=1.7cm,valign=top]
            \centering
            \begin{tikzpicture}\path (-0.5,-0.35) rectangle (0.5,0.35);\pic[insarouge,scale=0.9] at (0,0) {icone personne};\end{tikzpicture}\\[1pt]
            \gros{Duo, famille}\\[1pt]{\footnotesize 2 à 3 appareils}
          \end{carte}
        \end{column}
        \begin{column}{0.47\linewidth}
          \begin{carte}[height=1.7cm,valign=top]
            \centering
            \begin{tikzpicture}\path (-0.5,-0.35) rectangle (0.5,0.35);\pic[scale=0.9] at (0,0) {icone mont};\end{tikzpicture}\\[1pt]
            \gros{Randonnée}\\[1pt]{\footnotesize 4 appareils}
          \end{carte}
        \end{column}
      \end{columns}
      \vspace{0.1cm}
      \begin{columns}[T]
        \begin{column}{0.47\linewidth}
          \begin{carte}[height=1.7cm,valign=top]
            \centering
            \begin{tikzpicture}\path (-0.5,-0.35) rectangle (0.5,0.35);\pic[insarouge,scale=0.9] at (0,0) {icone tente};\end{tikzpicture}\\[1pt]
            \gros{Refuge, club}\\[1pt]{\footnotesize un relais et ses utilisateurs}
          \end{carte}
        \end{column}
        \begin{column}{0.47\linewidth}
          \begin{carte}[height=1.7cm,valign=top]
            \centering
            \begin{tikzpicture}\path (-0.5,-0.35) rectangle (0.5,0.35);\pic[scale=0.9] at (0,0) {icone croix};\end{tikzpicture}\\[1pt]
            \gros{ONG}\\[1pt]{\footnotesize secours~: 10 à 30 appareils}
          \end{carte}
        \end{column}
      \end{columns}
    \end{column}
  \end{columns}
  \vspace{0.05cm}
  \begin{pcompact}\centering\footnotesize
    Le bon produit est peut-être moins «\,une antenne\,» qu'un réseau prêt à l'emploi pour un groupe constitué.
  \end{pcompact}
  \srcnote{Micro-réseaux d'après le SWOT de l'équipe~; tailles de groupe à mesurer dans le questionnaire (questions 6 et 18).}
\end{frame}

%---------------------------------------------------------------------------
\begin{frame}{Enjeu 2~: un prix d'appel, une marge à protéger}
  \footnotesize\renewcommand{\arraystretch}{1.3}
  \hyphenpenalty=10000 \exhyphenpenalty=10000
  \newcolumntype{N}{>{\raggedleft\arraybackslash}p{1.5cm}}
  \setlength{\tabcolsep}{3pt}
  \noindent\begin{tabular}{@{}>{\raggedright\arraybackslash}p{2.9cm}NNNNNN@{}}
    \rowcolor{insarose}\textcolor{insarouge}{\bfseries Offre} & \textcolor{insarouge}{\bfseries Prix TTC} & \textcolor{insarouge}{\bfseries Par kit} & \textcolor{insarouge}{\bfseries Marge par commande} & \textcolor{insarouge}{\bfseries Marge par kit} & \textcolor{insarouge}{\bfseries Point mort (kits/an)} & \textcolor{insarouge}{\bfseries Résultat année 3}\\
    Kit seul (référence)    & 69\,€  & 69,00\,€ & 20,17\,€ & 20,17\,€ & 21\,300 & +13,7\,k€\\
    Solo                    & 79\,€  & 79,00\,€ & 28,50\,€ & 28,50\,€ & 15\,100 & +197\,k€\\
    \rowcolor{insabeige}Duo (2 kits)           & 139\,€ & 69,50\,€ & 45,17\,€ & 22,59\,€ & 19\,000 & +67\,k€\\
    Groupe (4 kits)         & 249\,€ & 62,25\,€ & 70,18\,€ & 17,54\,€ & 24\,500 & −44\,k€\\
  \end{tabular}
  \vspace{0.12cm}
  \begin{columns}[T]
    \begin{column}{0.325\textwidth}
      \begin{carte}[height=1.7cm,valign=top]
        \raggedright\gros{−73\,k€}\\[1pt]{\footnotesize sur le résultat de l'année 3 si le coût matériel monte de 10\,\% (+3\,€ par kit).}
      \end{carte}
    \end{column}
    \begin{column}{0.325\textwidth}
      \begin{carte}[height=1.7cm,valign=top]
        \raggedright\gros{−44\,k€}\\[1pt]{\footnotesize si la vente, le SAV et la logistique coûtent 2\,€ de plus par commande.}
      \end{carte}
    \end{column}
    \begin{column}{0.325\textwidth}
      \begin{carte}[height=1.7cm,valign=top]
        \raggedright\gros{±92\,k€}\\[1pt]{\footnotesize pour 5\,€ de prix en plus ou en moins sur chaque kit.}
      \end{carte}
    \end{column}
  \end{columns}
  \vspace{0.08cm}
  \begin{pcompact}\centering\footnotesize
    Le Duo relève la marge par kit~; le Groupe à 249\,€ la baisse (il faudrait ≥ 262\,€). Le prix se tranche par le test (questions 19 à 21 du questionnaire), pas par intuition.
  \end{pcompact}
  \srcnote{Calcul avec le modèle de la partie Finances~: matériel 30,30\,€, vente 4\,€ par commande, 1 kit offert pour 10, charges fixes 430\,k€, 22\,000 kits. Marge = prix HT moins coûts variables.}
\end{frame}

%---------------------------------------------------------------------------
\begin{frame}{Enjeu 3~: choisir un marché de départ}
  \begin{columns}[T]
    \begin{column}{0.35\textwidth}
      \begin{carte}[height=4.45cm,valign=top]
        \raggedright
        {\sffamily\bfseries\scriptsize\color{insarouge}CIBLE DE DÉPART (HYPOTHÈSE)}\\[2pt]
        {\sffamily\bfseries\large\color{insarouge}Randonneurs et groupes outdoor}\\[2pt]
        {\footnotesize\textbf{Pourquoi~?}}
        \begin{puces}
          \item besoin simple à comprendre
          \item usage réel en zone blanche
          \item marché large et accessible
          \item vente plus simple que le secours
        \end{puces}
        {\scriptsize\textcolor{insarouge}{À confirmer par l'étude quantitative.}}
      \end{carte}
    \end{column}
    \begin{column}{0.36\textwidth}
      {\sffamily\bfseries\footnotesize Progression naturelle}\\[3pt]
      \begin{tikzpicture}[font=\footnotesize]
        \draw[insarouge!40,thick] (0.2,0.05) -- (0.2,-3.55);
        \node[num] at (0.2,0)     {1};
        \node[num] at (0.2,-0.9) {2};
        \node[num] at (0.2,-1.8)  {3};
        \node[num] at (0.2,-2.7) {4};
        \node[num] at (0.2,-3.6)  {5};
        \node[anchor=west,text width=3.9cm,align=left] at (0.5,0)     {{\bfseries Outdoor}\\[-1pt]{\scriptsize randonneurs et groupes}};
        \node[anchor=west,text width=3.9cm,align=left] at (0.5,-0.9) {{\bfseries Refuges et clubs}\\[-1pt]{\scriptsize points de relais locaux}};
        \node[anchor=west,text width=3.9cm,align=left] at (0.5,-1.8)  {{\bfseries Associations}\\[-1pt]{\scriptsize ADRASEC, communautés structurées}};
        \node[anchor=west,text width=3.9cm,align=left] at (0.5,-2.7) {{\bfseries ONG, communes, pros}\\[-1pt]{\scriptsize déploiements plus larges}};
        \node[anchor=west,text width=3.9cm,align=left] at (0.5,-3.6)  {{\bfseries Grand public}\\[-1pt]{\scriptsize familles, voyageurs, une fois le réseau visible}};
      \end{tikzpicture}
    \end{column}
    \begin{column}{0.25\textwidth}
      \begin{idee}\footnotesize\centering
        {\sffamily\bfseries\color{insarouge}L'objectif}\\[3pt]
        Créer d'abord une zone où le produit est vraiment utile…\\[3pt]
        …puis utiliser cette densité, les retours terrain et la preuve d'usage pour élargir.
      \end{idee}
    \end{column}
  \end{columns}
  \vspace{0.05cm}
  \begin{pcompact}\centering\footnotesize
    Ne pas viser «\,tout le monde\,» au lancement~: gagner un premier territoire d'usage, puis élargir.
  \end{pcompact}
  \srcnote{Choix du marché de départ d'après le SWOT de l'équipe~; les seuils de décision sont dans la partie Étude quantitative.}
\end{frame}

%  Fichier de contenu : ni préambule ni \begin{document}, commence par \section.
%
%  Remplace l'ancien SWOT en PDF (\includepdf supprimé) : le diagnostic est maintenant écrit
%  en LaTeX/TikZ, avec la charte du reste de la présentation, et complété par un SWOT croisé.
%  Ce qui a changé par rapport à la version PDF :
%   - la décomposition des 69 € n'est plus répétée ici (elle est dans la partie Finances) ;
%   - chaque point du diagnostic renvoie à un « frein » F1 à F7 de l'étude documentaire ;
%   - les marges des offres Solo / Duo / Groupe sont recalculées (marge par commande ET par kit,
%     avec le coût de vente de 4 € compté une seule fois par commande) ;
%   - le SWOT croisé (stratégies SO, WO, ST, WT) relie le diagnostic aux trois enjeux.
%  Hypothèses chiffrées : celles de la partie Finances (kit 69 € TTC, matériel 30,30 €, vente 4 €
%  par commande, 1 kit offert pour 10, charges fixes de l'année 3 = 430 k€).
%  Les offres Solo (79 €), Duo (139 €, 2 kits) et Groupe (249 €, 4 kits) sont celles du SWOT d'origine ;
%  le nombre de kits du Groupe (4) est une hypothèse.
%%%%%%%%%%%%%%%%%%%%%%%%%%%%%%%%%%%%%%%%%%%%%%%%%%%%%%%%%%%%%%%%%%%%%%%%%%%%%%
\section{SWOT et enjeux}

% Entrée de SWOT : \swit[couleur du code]{code}{titre}{texte}  (titre et texte sur le même paragraphe)
\newcommand{\swit}[4][insarouge]{%
  \begin{tcolorbox}[colback=insabeige,colframe=insabeige,boxrule=0pt,arc=2pt,left=4pt,right=4pt,
                    top=1.5pt,bottom=1.5pt,before skip=0pt,after skip=2pt]
    \scriptsize\raggedright{\sffamily\bfseries\color{#1}#2}\hspace{0.3em}{\sffamily\bfseries\footnotesize #3}\hspace{0.3em}#4\par
  \end{tcolorbox}}
% En-tête de colonne : \swhead{Forces}{couleur}
\newcommand{\swhead}[2]{%
  {\sffamily\bfseries\footnotesize\color{#2}\MakeUppercase{#1}}\\[-3pt]{\color{#2}\rule{\linewidth}{0.6pt}}\\[-1pt]}
% renvoi vers un enjeu ou un frein
\newcommand{\swref}[1]{\textcolor{insarouge}{\textbf{→\,#1}}}

%---------------------------------------------------------------------------
\begin{frame}{Diagnostic interne~: forces et faiblesses}
  \begin{columns}[T]
    \begin{column}{0.485\textwidth}
      \swhead{Forces}{insarouge}
      \swit{S1}{Proposition de valeur claire}{Une messagerie texte sans réseau mobile, sans compte ni abonnement.}
      \swit{S2}{Coût d'entrée faible}{69\,€ une fois, contre ≈ 830\,€ sur trois ans pour une balise satellite.}
      \swit{S3}{Écosystème libre existant}{Meshtastic, Heltec V3, applications~: l'effort porte sur l'usage.}
      \swit{S4}{Simplicité comme différence}{Réglages Europe, clé privée, QR code, guides en français.}
      \swit{S5}{Un socle, plusieurs usages}{Outdoor, campagne, ONG, crise~: un kit, quatre terrains.}
    \end{column}
    \begin{column}{0.485\textwidth}
      \swhead{Faiblesses}{insagris}
      \swit[insagris]{W1}{Faible barrière technologique}{Un bricoleur refait un kit proche pour 18 à 30\,€. \swref{F5}}
      \swit[insagris]{W2}{Effet de réseau indispensable}{Un kit isolé ne sert à rien~: la valeur vient des voisins. \swref{enjeu 1}}
      \swit[insagris]{W3}{Marge de contribution fine}{20,17\,€ par kit, mais 0,62\,€ de résultat par kit en année 3. \swref{enjeu 2}}
      \swit[insagris]{W4}{Financement avant rentabilité}{≈ 450\,k€ à financer, point mort ≈ 21\,300 kits par an. \swref{F7}}
      \swit[insagris]{W5}{Positionnement trop large}{ONG, randonneurs, makers, public~: attentes différentes. \swref{enjeu 3}}
    \end{column}
  \end{columns}
  \vspace{-0.2cm}
  \begin{pcompact}\centering\footnotesize
    Lecture clé~: la technologie est accessible~; la valeur viendra de l'usage, de la marque et de l'écosystème.
  \end{pcompact}
  \srcnote{Synthèse des parties Concurrence, Finances et Étude documentaire. F2 à F7 = freins de l'étude documentaire.}
\end{frame}

%---------------------------------------------------------------------------
\begin{frame}{Diagnostic externe~: opportunités et menaces}
  \begin{columns}[T]
    \begin{column}{0.485\textwidth}
      \swhead{Opportunités}{insarouge}
      \swit{O1}{Un marché outdoor important}{27 millions de pratiquants, dont 240\,000 licenciés FFRandonnée.}
      \swit{O2}{Besoin croissant de résilience}{Zones blanches, crises~: l'État invite à préparer un kit de 72\,h.}
      \swit{O3}{Réseaux locaux à construire}{Refuges, clubs, communes, ADRASEC~: points d'ancrage du maillage.}
      \swit{O4}{Précommandes adaptées}{Le financement participatif teste la demande et paie le premier stock.}
      \swit{O5}{Made in France, souveraineté}{Relocalisation par étapes~: un argument pour les collectivités.}
    \end{column}
    \begin{column}{0.485\textwidth}
      \swhead{Menaces}{insagris}
      \swit[insagris]{T1}{Fait-maison, kits concurrents}{Cartes dès 18\,€, kits prêts à l'emploi 59 à 153\,€. \swref{C}}
      \swit[insagris]{T2}{Substituts puissants}{SOS satellite des smartphones, balises, SMS satellite. \swref{F2}}
      \swit[insagris]{T3}{Contraintes réglementaires}{Marquage CE, cybersécurité, responsabilité produit. \swref{F3, F4}}
      \swit[insagris]{T4}{Dépendance amont}{Un seul fabricant de carte~; cap et marque de Meshtastic. \swref{F5, F6}}
      \swit[insagris]{T5}{Promesse mal comprise}{Ni téléphone, ni balise de détresse~: cadrer la promesse. \swref{F3}}
    \end{column}
  \end{columns}
  \vspace{-0.22cm}
  \begin{pcompact}\centering\footnotesize
    Opportunité réelle, mais fenêtre concurrentielle courte~: il faut bâtir vite une communauté et une marque.
  \end{pcompact}
  \srcnote{Sources~: \cite{ffrando,ffrando_licencies,guide_tous_resp}. F2 à F6 = freins de l'étude documentaire~; C = chapitre Concurrence.}
\end{frame}

%---------------------------------------------------------------------------
\newenvironment{pucesx}
  {\begingroup\setbeamerfont{itemize/enumerate body}{size=\scriptsize}\begin{itemize}\setlength{\itemsep}{0pt}\setlength{\topsep}{1pt}}
  {\end{itemize}\endgroup}
\begin{frame}{SWOT croisé~: quelles stratégies en déduire~?}
  \footnotesize\renewcommand{\arraystretch}{1.0}\setlength{\tabcolsep}{4pt}
  \hyphenpenalty=10000 \exhyphenpenalty=10000
  \noindent\begin{tabularx}{\textwidth}{@{}>{\raggedright\arraybackslash}p{1.65cm}>{\raggedright\arraybackslash}X>{\raggedright\arraybackslash}X@{}}
    & \cellcolor{insarose}\textcolor{insarouge}{\bfseries Opportunités (O)} & \cellcolor{insarose}\textcolor{insarouge}{\bfseries Menaces (T)} \\[2pt]
    \cellcolor{insarose}\textcolor{insarouge}{\bfseries Forces (S)} &
      \cellcolor{insabeige}{\sffamily\bfseries\color{insarouge}Attaquer}
      \begin{pucesx}
        \item Micro-réseaux prêts à l'emploi (Duo, Groupe) pour clubs, refuges, groupes de randonnée.
        \item Précommandes~: tester la demande et financer le premier stock.
        \item Pilotes avec des ADRASEC et des communes~: preuves d'usage, premiers nœuds.
      \end{pucesx} &
      \cellcolor{insabeige}{\sffamily\bfseries\color{insarouge}Défendre}
      \begin{pucesx}
        \item Vendre le prêt-à-l'emploi~: réglages, clé privée, guide, garantie, marquage CE.
        \item Rester un complément du smartphone et de la balise, sans promesse de secours.
        \item Contribuer au code libre, rester « compatible Meshtastic ».
      \end{pucesx} \\[3pt]
    \cellcolor{insarose}\textcolor{insarouge}{\bfseries Faiblesses (W)} &
      \cellcolor{insabeige}{\sffamily\bfseries\color{insarouge}Corriger}
      \begin{pucesx}
        \item Densifier autour de points d'ancrage (refuge, club, mairie) contre l'effet de réseau.
        \item Tester Solo, Duo et Groupe avant la série pour relever la marge.
        \item Bâtir le dossier de financement sur les précommandes.
      \end{pucesx} &
      \cellcolor{insabeige}{\sffamily\bfseries\color{insarouge}Éviter}
      \begin{pucesx}
        \item Notice claire, test à 100\,\%, assurance responsabilité produit.
        \item Double source de cartes (V3 et V4), stock tampon, vente d'abord en France.
        \item Fixer des seuils d'arrêt avant le moule et la série (étude quantitative).
      \end{pucesx} \\
  \end{tabularx}
  \srcnote{Croisement des deux diapositives précédentes (méthode TOWS)~: S, W = forces et faiblesses~; O, T = opportunités et menaces.}
\end{frame}

%---------------------------------------------------------------------------
% Liaisons possibles entre n kits : n(n-1)/2 (2 kits = 1, 4 = 6, 6 = 15, 10 = 45)
\begin{frame}{Enjeu 1~: résoudre l'effet de réseau}
  \begin{columns}[T]
    \begin{column}{0.46\textwidth}
      \centering
      \begin{tikzpicture}[font=\footnotesize]
        \node[pcli,text width=2.3cm,minimum height=0.7cm] (a) at (0,1.05) {Peu d'utilisateurs};
        \node[pgris,text width=1.6cm,minimum height=0.7cm] (b) at (3.55,1.05) {Peu de relais};
        \node[pgris,text width=1.6cm,minimum height=0.7cm] (c) at (1.78,-0.3) {Valeur limitée};
        \draw[fluxr] (a) -- (b);
        \draw[fluxr] (b) -- (c);
        \draw[fluxr] (c) -- (a);
        \node[font=\scriptsize\bfseries,text=insarouge,align=center] at (1.78,0.5) {cercle\\vicieux};
      \end{tikzpicture}

      \vspace{0.05cm}
      {\footnotesize\textbf{Un groupe vaut plus que la somme de ses kits}~: $n$ kits offrent $n(n{-}1)/2$ liaisons.\par}
      \begin{tikzpicture}[font=\scriptsize]
        \tikzset{noeud/.style={circle,fill=insarouge,inner sep=1.3pt}}
        % 2 kits
        \begin{scope}[shift={(0.0,0)}]
          \node[noeud] (u0) at (-0.28,0) {};\node[noeud] (u1) at (0.28,0) {};
          \draw[insagris,thick] (u0) -- (u1);
          \node[align=center,anchor=north] at (0,-0.55) {2 kits\\\textbf{1 liaison}};
        \end{scope}
        % 4 kits
        \begin{scope}[shift={(1.85,0)}]
          \foreach \i in {0,...,3} {\node[noeud] (v\i) at ({45+90*\i}:0.4) {};}
          \foreach \i/\j in {0/1,0/2,0/3,1/2,1/3,2/3} {\draw[insagris,thick] (v\i) -- (v\j);}
          \node[align=center,anchor=north] at (0,-0.55) {4 kits\\\textbf{6 liaisons}};
        \end{scope}
        % 6 kits
        \begin{scope}[shift={(3.7,0)}]
          \foreach \i in {0,...,5} {\node[noeud] (w\i) at ({90+60*\i}:0.42) {};}
          \foreach \i/\j in {0/1,0/2,0/3,0/4,0/5,1/2,1/3,1/4,1/5,2/3,2/4,2/5,3/4,3/5,4/5} {\draw[insagris,thin] (w\i) -- (w\j);}
          \node[align=center,anchor=north] at (0,-0.55) {6 kits\\\textbf{15 liaisons}};
        \end{scope}
        \node[align=center,font=\scriptsize] at (5.15,0) {10 kits\\\textbf{45}};
      \end{tikzpicture}
    \end{column}
    \begin{column}{0.51\textwidth}
      \begin{columns}[T]
        \begin{column}{0.47\linewidth}
          \begin{carte}[height=1.7cm,valign=top]
            \centering
            \begin{tikzpicture}\path (-0.5,-0.35) rectangle (0.5,0.35);\pic[insarouge,scale=0.9] at (0,0) {icone personne};\end{tikzpicture}\\[1pt]
            \gros{Duo, famille}\\[1pt]{\footnotesize 2 à 3 appareils}
          \end{carte}
        \end{column}
        \begin{column}{0.47\linewidth}
          \begin{carte}[height=1.7cm,valign=top]
            \centering
            \begin{tikzpicture}\path (-0.5,-0.35) rectangle (0.5,0.35);\pic[scale=0.9] at (0,0) {icone mont};\end{tikzpicture}\\[1pt]
            \gros{Randonnée}\\[1pt]{\footnotesize 4 appareils}
          \end{carte}
        \end{column}
      \end{columns}
      \vspace{0.1cm}
      \begin{columns}[T]
        \begin{column}{0.47\linewidth}
          \begin{carte}[height=1.7cm,valign=top]
            \centering
            \begin{tikzpicture}\path (-0.5,-0.35) rectangle (0.5,0.35);\pic[insarouge,scale=0.9] at (0,0) {icone tente};\end{tikzpicture}\\[1pt]
            \gros{Refuge, club}\\[1pt]{\footnotesize un relais et ses utilisateurs}
          \end{carte}
        \end{column}
        \begin{column}{0.47\linewidth}
          \begin{carte}[height=1.7cm,valign=top]
            \centering
            \begin{tikzpicture}\path (-0.5,-0.35) rectangle (0.5,0.35);\pic[scale=0.9] at (0,0) {icone croix};\end{tikzpicture}\\[1pt]
            \gros{ONG}\\[1pt]{\footnotesize secours~: 10 à 30 appareils}
          \end{carte}
        \end{column}
      \end{columns}
    \end{column}
  \end{columns}
  \vspace{0.05cm}
  \begin{pcompact}\centering\footnotesize
    Le bon produit est peut-être moins «\,une antenne\,» qu'un réseau prêt à l'emploi pour un groupe constitué.
  \end{pcompact}
  \srcnote{Micro-réseaux d'après le SWOT de l'équipe~; tailles de groupe à mesurer dans le questionnaire (questions 6 et 18).}
\end{frame}

%---------------------------------------------------------------------------
\begin{frame}{Enjeu 2~: un prix d'appel, une marge à protéger}
  \footnotesize\renewcommand{\arraystretch}{1.3}
  \hyphenpenalty=10000 \exhyphenpenalty=10000
  \newcolumntype{N}{>{\raggedleft\arraybackslash}p{1.5cm}}
  \setlength{\tabcolsep}{3pt}
  \noindent\begin{tabular}{@{}>{\raggedright\arraybackslash}p{2.9cm}NNNNNN@{}}
    \rowcolor{insarose}\textcolor{insarouge}{\bfseries Offre} & \textcolor{insarouge}{\bfseries Prix TTC} & \textcolor{insarouge}{\bfseries Par kit} & \textcolor{insarouge}{\bfseries Marge par commande} & \textcolor{insarouge}{\bfseries Marge par kit} & \textcolor{insarouge}{\bfseries Point mort (kits/an)} & \textcolor{insarouge}{\bfseries Résultat année 3}\\
    Kit seul (référence)    & 69\,€  & 69,00\,€ & 20,17\,€ & 20,17\,€ & 21\,300 & +13,7\,k€\\
    Solo                    & 79\,€  & 79,00\,€ & 28,50\,€ & 28,50\,€ & 15\,100 & +197\,k€\\
    \rowcolor{insabeige}Duo (2 kits)           & 139\,€ & 69,50\,€ & 45,17\,€ & 22,59\,€ & 19\,000 & +67\,k€\\
    Groupe (4 kits)         & 249\,€ & 62,25\,€ & 70,18\,€ & 17,54\,€ & 24\,500 & −44\,k€\\
  \end{tabular}
  \vspace{0.12cm}
  \begin{columns}[T]
    \begin{column}{0.325\textwidth}
      \begin{carte}[height=1.7cm,valign=top]
        \raggedright\gros{−73\,k€}\\[1pt]{\footnotesize sur le résultat de l'année 3 si le coût matériel monte de 10\,\% (+3\,€ par kit).}
      \end{carte}
    \end{column}
    \begin{column}{0.325\textwidth}
      \begin{carte}[height=1.7cm,valign=top]
        \raggedright\gros{−44\,k€}\\[1pt]{\footnotesize si la vente, le SAV et la logistique coûtent 2\,€ de plus par commande.}
      \end{carte}
    \end{column}
    \begin{column}{0.325\textwidth}
      \begin{carte}[height=1.7cm,valign=top]
        \raggedright\gros{±92\,k€}\\[1pt]{\footnotesize pour 5\,€ de prix en plus ou en moins sur chaque kit.}
      \end{carte}
    \end{column}
  \end{columns}
  \vspace{0.08cm}
  \begin{pcompact}\centering\footnotesize
    Le Duo relève la marge par kit~; le Groupe à 249\,€ la baisse (il faudrait ≥ 262\,€). Le prix se tranche par le test (questions 19 à 21 du questionnaire), pas par intuition.
  \end{pcompact}
  \srcnote{Calcul avec le modèle de la partie Finances~: matériel 30,30\,€, vente 4\,€ par commande, 1 kit offert pour 10, charges fixes 430\,k€, 22\,000 kits. Marge = prix HT moins coûts variables.}
\end{frame}

%---------------------------------------------------------------------------
\begin{frame}{Enjeu 3~: choisir un marché de départ}
  \begin{columns}[T]
    \begin{column}{0.35\textwidth}
      \begin{carte}[height=4.45cm,valign=top]
        \raggedright
        {\sffamily\bfseries\scriptsize\color{insarouge}CIBLE DE DÉPART (HYPOTHÈSE)}\\[2pt]
        {\sffamily\bfseries\large\color{insarouge}Randonneurs et groupes outdoor}\\[2pt]
        {\footnotesize\textbf{Pourquoi~?}}
        \begin{puces}
          \item besoin simple à comprendre
          \item usage réel en zone blanche
          \item marché large et accessible
          \item vente plus simple que le secours
        \end{puces}
        {\scriptsize\textcolor{insarouge}{À confirmer par l'étude quantitative.}}
      \end{carte}
    \end{column}
    \begin{column}{0.36\textwidth}
      {\sffamily\bfseries\footnotesize Progression naturelle}\\[3pt]
      \begin{tikzpicture}[font=\footnotesize]
        \draw[insarouge!40,thick] (0.2,0.05) -- (0.2,-3.55);
        \node[num] at (0.2,0)     {1};
        \node[num] at (0.2,-0.9) {2};
        \node[num] at (0.2,-1.8)  {3};
        \node[num] at (0.2,-2.7) {4};
        \node[num] at (0.2,-3.6)  {5};
        \node[anchor=west,text width=3.9cm,align=left] at (0.5,0)     {{\bfseries Outdoor}\\[-1pt]{\scriptsize randonneurs et groupes}};
        \node[anchor=west,text width=3.9cm,align=left] at (0.5,-0.9) {{\bfseries Refuges et clubs}\\[-1pt]{\scriptsize points de relais locaux}};
        \node[anchor=west,text width=3.9cm,align=left] at (0.5,-1.8)  {{\bfseries Associations}\\[-1pt]{\scriptsize ADRASEC, communautés structurées}};
        \node[anchor=west,text width=3.9cm,align=left] at (0.5,-2.7) {{\bfseries ONG, communes, pros}\\[-1pt]{\scriptsize déploiements plus larges}};
        \node[anchor=west,text width=3.9cm,align=left] at (0.5,-3.6)  {{\bfseries Grand public}\\[-1pt]{\scriptsize familles, voyageurs, une fois le réseau visible}};
      \end{tikzpicture}
    \end{column}
    \begin{column}{0.25\textwidth}
      \begin{idee}\footnotesize\centering
        {\sffamily\bfseries\color{insarouge}L'objectif}\\[3pt]
        Créer d'abord une zone où le produit est vraiment utile…\\[3pt]
        …puis utiliser cette densité, les retours terrain et la preuve d'usage pour élargir.
      \end{idee}
    \end{column}
  \end{columns}
  \vspace{0.05cm}
  \begin{pcompact}\centering\footnotesize
    Ne pas viser «\,tout le monde\,» au lancement~: gagner un premier territoire d'usage, puis élargir.
  \end{pcompact}
  \srcnote{Choix du marché de départ d'après le SWOT de l'équipe~; les seuils de décision sont dans la partie Étude quantitative.}
\end{frame}

%  Fichier de contenu : ni préambule ni \begin{document}, commence par \section.
%
%  Remplace l'ancien SWOT en PDF (\includepdf supprimé) : le diagnostic est maintenant écrit
%  en LaTeX/TikZ, avec la charte du reste de la présentation, et complété par un SWOT croisé.
%  Ce qui a changé par rapport à la version PDF :
%   - la décomposition des 69 € n'est plus répétée ici (elle est dans la partie Finances) ;
%   - chaque point du diagnostic renvoie à un « frein » F1 à F7 de l'étude documentaire ;
%   - les marges des offres Solo / Duo / Groupe sont recalculées (marge par commande ET par kit,
%     avec le coût de vente de 4 € compté une seule fois par commande) ;
%   - le SWOT croisé (stratégies SO, WO, ST, WT) relie le diagnostic aux trois enjeux.
%  Hypothèses chiffrées : celles de la partie Finances (kit 69 € TTC, matériel 30,30 €, vente 4 €
%  par commande, 1 kit offert pour 10, charges fixes de l'année 3 = 430 k€).
%  Les offres Solo (79 €), Duo (139 €, 2 kits) et Groupe (249 €, 4 kits) sont celles du SWOT d'origine ;
%  le nombre de kits du Groupe (4) est une hypothèse.
%%%%%%%%%%%%%%%%%%%%%%%%%%%%%%%%%%%%%%%%%%%%%%%%%%%%%%%%%%%%%%%%%%%%%%%%%%%%%%
\section{SWOT et enjeux}

% Entrée de SWOT : \swit[couleur du code]{code}{titre}{texte}  (titre et texte sur le même paragraphe)
\newcommand{\swit}[4][insarouge]{%
  \begin{tcolorbox}[colback=insabeige,colframe=insabeige,boxrule=0pt,arc=2pt,left=4pt,right=4pt,
                    top=1.5pt,bottom=1.5pt,before skip=0pt,after skip=2pt]
    \scriptsize\raggedright{\sffamily\bfseries\color{#1}#2}\hspace{0.3em}{\sffamily\bfseries\footnotesize #3}\hspace{0.3em}#4\par
  \end{tcolorbox}}
% En-tête de colonne : \swhead{Forces}{couleur}
\newcommand{\swhead}[2]{%
  {\sffamily\bfseries\footnotesize\color{#2}\MakeUppercase{#1}}\\[-3pt]{\color{#2}\rule{\linewidth}{0.6pt}}\\[-1pt]}
% renvoi vers un enjeu ou un frein
\newcommand{\swref}[1]{\textcolor{insarouge}{\textbf{→\,#1}}}

%---------------------------------------------------------------------------
\begin{frame}{Diagnostic interne~: forces et faiblesses}
  \begin{columns}[T]
    \begin{column}{0.485\textwidth}
      \swhead{Forces}{insarouge}
      \swit{S1}{Proposition de valeur claire}{Une messagerie texte sans réseau mobile, sans compte ni abonnement.}
      \swit{S2}{Coût d'entrée faible}{69\,€ une fois, contre ≈ 830\,€ sur trois ans pour une balise satellite.}
      \swit{S3}{Écosystème libre existant}{Meshtastic, Heltec V3, applications~: l'effort porte sur l'usage.}
      \swit{S4}{Simplicité comme différence}{Réglages Europe, clé privée, QR code, guides en français.}
      \swit{S5}{Un socle, plusieurs usages}{Outdoor, campagne, ONG, crise~: un kit, quatre terrains.}
    \end{column}
    \begin{column}{0.485\textwidth}
      \swhead{Faiblesses}{insagris}
      \swit[insagris]{W1}{Faible barrière technologique}{Un bricoleur refait un kit proche pour 18 à 30\,€. \swref{F5}}
      \swit[insagris]{W2}{Effet de réseau indispensable}{Un kit isolé ne sert à rien~: la valeur vient des voisins. \swref{enjeu 1}}
      \swit[insagris]{W3}{Marge de contribution fine}{20,17\,€ par kit, mais 0,62\,€ de résultat par kit en année 3. \swref{enjeu 2}}
      \swit[insagris]{W4}{Financement avant rentabilité}{≈ 450\,k€ à financer, point mort ≈ 21\,300 kits par an. \swref{F7}}
      \swit[insagris]{W5}{Positionnement trop large}{ONG, randonneurs, makers, public~: attentes différentes. \swref{enjeu 3}}
    \end{column}
  \end{columns}
  \vspace{-0.2cm}
  \begin{pcompact}\centering\footnotesize
    Lecture clé~: la technologie est accessible~; la valeur viendra de l'usage, de la marque et de l'écosystème.
  \end{pcompact}
  \srcnote{Synthèse des parties Concurrence, Finances et Étude documentaire. F2 à F7 = freins de l'étude documentaire.}
\end{frame}

%---------------------------------------------------------------------------
\begin{frame}{Diagnostic externe~: opportunités et menaces}
  \begin{columns}[T]
    \begin{column}{0.485\textwidth}
      \swhead{Opportunités}{insarouge}
      \swit{O1}{Un marché outdoor important}{27 millions de pratiquants, dont 240\,000 licenciés FFRandonnée.}
      \swit{O2}{Besoin croissant de résilience}{Zones blanches, crises~: l'État invite à préparer un kit de 72\,h.}
      \swit{O3}{Réseaux locaux à construire}{Refuges, clubs, communes, ADRASEC~: points d'ancrage du maillage.}
      \swit{O4}{Précommandes adaptées}{Le financement participatif teste la demande et paie le premier stock.}
      \swit{O5}{Made in France, souveraineté}{Relocalisation par étapes~: un argument pour les collectivités.}
    \end{column}
    \begin{column}{0.485\textwidth}
      \swhead{Menaces}{insagris}
      \swit[insagris]{T1}{Fait-maison, kits concurrents}{Cartes dès 18\,€, kits prêts à l'emploi 59 à 153\,€. \swref{C}}
      \swit[insagris]{T2}{Substituts puissants}{SOS satellite des smartphones, balises, SMS satellite. \swref{F2}}
      \swit[insagris]{T3}{Contraintes réglementaires}{Marquage CE, cybersécurité, responsabilité produit. \swref{F3, F4}}
      \swit[insagris]{T4}{Dépendance amont}{Un seul fabricant de carte~; cap et marque de Meshtastic. \swref{F5, F6}}
      \swit[insagris]{T5}{Promesse mal comprise}{Ni téléphone, ni balise de détresse~: cadrer la promesse. \swref{F3}}
    \end{column}
  \end{columns}
  \vspace{-0.22cm}
  \begin{pcompact}\centering\footnotesize
    Opportunité réelle, mais fenêtre concurrentielle courte~: il faut bâtir vite une communauté et une marque.
  \end{pcompact}
  \srcnote{Sources~: \cite{ffrando,ffrando_licencies,guide_tous_resp}. F2 à F6 = freins de l'étude documentaire~; C = chapitre Concurrence.}
\end{frame}

%---------------------------------------------------------------------------
\newenvironment{pucesx}
  {\begingroup\setbeamerfont{itemize/enumerate body}{size=\scriptsize}\begin{itemize}\setlength{\itemsep}{0pt}\setlength{\topsep}{1pt}}
  {\end{itemize}\endgroup}
\begin{frame}{SWOT croisé~: quelles stratégies en déduire~?}
  \footnotesize\renewcommand{\arraystretch}{1.0}\setlength{\tabcolsep}{4pt}
  \hyphenpenalty=10000 \exhyphenpenalty=10000
  \noindent\begin{tabularx}{\textwidth}{@{}>{\raggedright\arraybackslash}p{1.65cm}>{\raggedright\arraybackslash}X>{\raggedright\arraybackslash}X@{}}
    & \cellcolor{insarose}\textcolor{insarouge}{\bfseries Opportunités (O)} & \cellcolor{insarose}\textcolor{insarouge}{\bfseries Menaces (T)} \\[2pt]
    \cellcolor{insarose}\textcolor{insarouge}{\bfseries Forces (S)} &
      \cellcolor{insabeige}{\sffamily\bfseries\color{insarouge}Attaquer}
      \begin{pucesx}
        \item Micro-réseaux prêts à l'emploi (Duo, Groupe) pour clubs, refuges, groupes de randonnée.
        \item Précommandes~: tester la demande et financer le premier stock.
        \item Pilotes avec des ADRASEC et des communes~: preuves d'usage, premiers nœuds.
      \end{pucesx} &
      \cellcolor{insabeige}{\sffamily\bfseries\color{insarouge}Défendre}
      \begin{pucesx}
        \item Vendre le prêt-à-l'emploi~: réglages, clé privée, guide, garantie, marquage CE.
        \item Rester un complément du smartphone et de la balise, sans promesse de secours.
        \item Contribuer au code libre, rester « compatible Meshtastic ».
      \end{pucesx} \\[3pt]
    \cellcolor{insarose}\textcolor{insarouge}{\bfseries Faiblesses (W)} &
      \cellcolor{insabeige}{\sffamily\bfseries\color{insarouge}Corriger}
      \begin{pucesx}
        \item Densifier autour de points d'ancrage (refuge, club, mairie) contre l'effet de réseau.
        \item Tester Solo, Duo et Groupe avant la série pour relever la marge.
        \item Bâtir le dossier de financement sur les précommandes.
      \end{pucesx} &
      \cellcolor{insabeige}{\sffamily\bfseries\color{insarouge}Éviter}
      \begin{pucesx}
        \item Notice claire, test à 100\,\%, assurance responsabilité produit.
        \item Double source de cartes (V3 et V4), stock tampon, vente d'abord en France.
        \item Fixer des seuils d'arrêt avant le moule et la série (étude quantitative).
      \end{pucesx} \\
  \end{tabularx}
  \srcnote{Croisement des deux diapositives précédentes (méthode TOWS)~: S, W = forces et faiblesses~; O, T = opportunités et menaces.}
\end{frame}

%---------------------------------------------------------------------------
% Liaisons possibles entre n kits : n(n-1)/2 (2 kits = 1, 4 = 6, 6 = 15, 10 = 45)
\begin{frame}{Enjeu 1~: résoudre l'effet de réseau}
  \begin{columns}[T]
    \begin{column}{0.46\textwidth}
      \centering
      \begin{tikzpicture}[font=\footnotesize]
        \node[pcli,text width=2.3cm,minimum height=0.7cm] (a) at (0,1.05) {Peu d'utilisateurs};
        \node[pgris,text width=1.6cm,minimum height=0.7cm] (b) at (3.55,1.05) {Peu de relais};
        \node[pgris,text width=1.6cm,minimum height=0.7cm] (c) at (1.78,-0.3) {Valeur limitée};
        \draw[fluxr] (a) -- (b);
        \draw[fluxr] (b) -- (c);
        \draw[fluxr] (c) -- (a);
        \node[font=\scriptsize\bfseries,text=insarouge,align=center] at (1.78,0.5) {cercle\\vicieux};
      \end{tikzpicture}

      \vspace{0.05cm}
      {\footnotesize\textbf{Un groupe vaut plus que la somme de ses kits}~: $n$ kits offrent $n(n{-}1)/2$ liaisons.\par}
      \begin{tikzpicture}[font=\scriptsize]
        \tikzset{noeud/.style={circle,fill=insarouge,inner sep=1.3pt}}
        % 2 kits
        \begin{scope}[shift={(0.0,0)}]
          \node[noeud] (u0) at (-0.28,0) {};\node[noeud] (u1) at (0.28,0) {};
          \draw[insagris,thick] (u0) -- (u1);
          \node[align=center,anchor=north] at (0,-0.55) {2 kits\\\textbf{1 liaison}};
        \end{scope}
        % 4 kits
        \begin{scope}[shift={(1.85,0)}]
          \foreach \i in {0,...,3} {\node[noeud] (v\i) at ({45+90*\i}:0.4) {};}
          \foreach \i/\j in {0/1,0/2,0/3,1/2,1/3,2/3} {\draw[insagris,thick] (v\i) -- (v\j);}
          \node[align=center,anchor=north] at (0,-0.55) {4 kits\\\textbf{6 liaisons}};
        \end{scope}
        % 6 kits
        \begin{scope}[shift={(3.7,0)}]
          \foreach \i in {0,...,5} {\node[noeud] (w\i) at ({90+60*\i}:0.42) {};}
          \foreach \i/\j in {0/1,0/2,0/3,0/4,0/5,1/2,1/3,1/4,1/5,2/3,2/4,2/5,3/4,3/5,4/5} {\draw[insagris,thin] (w\i) -- (w\j);}
          \node[align=center,anchor=north] at (0,-0.55) {6 kits\\\textbf{15 liaisons}};
        \end{scope}
        \node[align=center,font=\scriptsize] at (5.15,0) {10 kits\\\textbf{45}};
      \end{tikzpicture}
    \end{column}
    \begin{column}{0.51\textwidth}
      \begin{columns}[T]
        \begin{column}{0.47\linewidth}
          \begin{carte}[height=1.7cm,valign=top]
            \centering
            \begin{tikzpicture}\path (-0.5,-0.35) rectangle (0.5,0.35);\pic[insarouge,scale=0.9] at (0,0) {icone personne};\end{tikzpicture}\\[1pt]
            \gros{Duo, famille}\\[1pt]{\footnotesize 2 à 3 appareils}
          \end{carte}
        \end{column}
        \begin{column}{0.47\linewidth}
          \begin{carte}[height=1.7cm,valign=top]
            \centering
            \begin{tikzpicture}\path (-0.5,-0.35) rectangle (0.5,0.35);\pic[scale=0.9] at (0,0) {icone mont};\end{tikzpicture}\\[1pt]
            \gros{Randonnée}\\[1pt]{\footnotesize 4 appareils}
          \end{carte}
        \end{column}
      \end{columns}
      \vspace{0.1cm}
      \begin{columns}[T]
        \begin{column}{0.47\linewidth}
          \begin{carte}[height=1.7cm,valign=top]
            \centering
            \begin{tikzpicture}\path (-0.5,-0.35) rectangle (0.5,0.35);\pic[insarouge,scale=0.9] at (0,0) {icone tente};\end{tikzpicture}\\[1pt]
            \gros{Refuge, club}\\[1pt]{\footnotesize un relais et ses utilisateurs}
          \end{carte}
        \end{column}
        \begin{column}{0.47\linewidth}
          \begin{carte}[height=1.7cm,valign=top]
            \centering
            \begin{tikzpicture}\path (-0.5,-0.35) rectangle (0.5,0.35);\pic[scale=0.9] at (0,0) {icone croix};\end{tikzpicture}\\[1pt]
            \gros{ONG}\\[1pt]{\footnotesize secours~: 10 à 30 appareils}
          \end{carte}
        \end{column}
      \end{columns}
    \end{column}
  \end{columns}
  \vspace{0.05cm}
  \begin{pcompact}\centering\footnotesize
    Le bon produit est peut-être moins «\,une antenne\,» qu'un réseau prêt à l'emploi pour un groupe constitué.
  \end{pcompact}
  \srcnote{Micro-réseaux d'après le SWOT de l'équipe~; tailles de groupe à mesurer dans le questionnaire (questions 6 et 18).}
\end{frame}

%---------------------------------------------------------------------------
\begin{frame}{Enjeu 2~: un prix d'appel, une marge à protéger}
  \footnotesize\renewcommand{\arraystretch}{1.3}
  \hyphenpenalty=10000 \exhyphenpenalty=10000
  \newcolumntype{N}{>{\raggedleft\arraybackslash}p{1.5cm}}
  \setlength{\tabcolsep}{3pt}
  \noindent\begin{tabular}{@{}>{\raggedright\arraybackslash}p{2.9cm}NNNNNN@{}}
    \rowcolor{insarose}\textcolor{insarouge}{\bfseries Offre} & \textcolor{insarouge}{\bfseries Prix TTC} & \textcolor{insarouge}{\bfseries Par kit} & \textcolor{insarouge}{\bfseries Marge par commande} & \textcolor{insarouge}{\bfseries Marge par kit} & \textcolor{insarouge}{\bfseries Point mort (kits/an)} & \textcolor{insarouge}{\bfseries Résultat année 3}\\
    Kit seul (référence)    & 69\,€  & 69,00\,€ & 20,17\,€ & 20,17\,€ & 21\,300 & +13,7\,k€\\
    Solo                    & 79\,€  & 79,00\,€ & 28,50\,€ & 28,50\,€ & 15\,100 & +197\,k€\\
    \rowcolor{insabeige}Duo (2 kits)           & 139\,€ & 69,50\,€ & 45,17\,€ & 22,59\,€ & 19\,000 & +67\,k€\\
    Groupe (4 kits)         & 249\,€ & 62,25\,€ & 70,18\,€ & 17,54\,€ & 24\,500 & −44\,k€\\
  \end{tabular}
  \vspace{0.12cm}
  \begin{columns}[T]
    \begin{column}{0.325\textwidth}
      \begin{carte}[height=1.7cm,valign=top]
        \raggedright\gros{−73\,k€}\\[1pt]{\footnotesize sur le résultat de l'année 3 si le coût matériel monte de 10\,\% (+3\,€ par kit).}
      \end{carte}
    \end{column}
    \begin{column}{0.325\textwidth}
      \begin{carte}[height=1.7cm,valign=top]
        \raggedright\gros{−44\,k€}\\[1pt]{\footnotesize si la vente, le SAV et la logistique coûtent 2\,€ de plus par commande.}
      \end{carte}
    \end{column}
    \begin{column}{0.325\textwidth}
      \begin{carte}[height=1.7cm,valign=top]
        \raggedright\gros{±92\,k€}\\[1pt]{\footnotesize pour 5\,€ de prix en plus ou en moins sur chaque kit.}
      \end{carte}
    \end{column}
  \end{columns}
  \vspace{0.08cm}
  \begin{pcompact}\centering\footnotesize
    Le Duo relève la marge par kit~; le Groupe à 249\,€ la baisse (il faudrait ≥ 262\,€). Le prix se tranche par le test (questions 19 à 21 du questionnaire), pas par intuition.
  \end{pcompact}
  \srcnote{Calcul avec le modèle de la partie Finances~: matériel 30,30\,€, vente 4\,€ par commande, 1 kit offert pour 10, charges fixes 430\,k€, 22\,000 kits. Marge = prix HT moins coûts variables.}
\end{frame}

%---------------------------------------------------------------------------
\begin{frame}{Enjeu 3~: choisir un marché de départ}
  \begin{columns}[T]
    \begin{column}{0.35\textwidth}
      \begin{carte}[height=4.45cm,valign=top]
        \raggedright
        {\sffamily\bfseries\scriptsize\color{insarouge}CIBLE DE DÉPART (HYPOTHÈSE)}\\[2pt]
        {\sffamily\bfseries\large\color{insarouge}Randonneurs et groupes outdoor}\\[2pt]
        {\footnotesize\textbf{Pourquoi~?}}
        \begin{puces}
          \item besoin simple à comprendre
          \item usage réel en zone blanche
          \item marché large et accessible
          \item vente plus simple que le secours
        \end{puces}
        {\scriptsize\textcolor{insarouge}{À confirmer par l'étude quantitative.}}
      \end{carte}
    \end{column}
    \begin{column}{0.36\textwidth}
      {\sffamily\bfseries\footnotesize Progression naturelle}\\[3pt]
      \begin{tikzpicture}[font=\footnotesize]
        \draw[insarouge!40,thick] (0.2,0.05) -- (0.2,-3.55);
        \node[num] at (0.2,0)     {1};
        \node[num] at (0.2,-0.9) {2};
        \node[num] at (0.2,-1.8)  {3};
        \node[num] at (0.2,-2.7) {4};
        \node[num] at (0.2,-3.6)  {5};
        \node[anchor=west,text width=3.9cm,align=left] at (0.5,0)     {{\bfseries Outdoor}\\[-1pt]{\scriptsize randonneurs et groupes}};
        \node[anchor=west,text width=3.9cm,align=left] at (0.5,-0.9) {{\bfseries Refuges et clubs}\\[-1pt]{\scriptsize points de relais locaux}};
        \node[anchor=west,text width=3.9cm,align=left] at (0.5,-1.8)  {{\bfseries Associations}\\[-1pt]{\scriptsize ADRASEC, communautés structurées}};
        \node[anchor=west,text width=3.9cm,align=left] at (0.5,-2.7) {{\bfseries ONG, communes, pros}\\[-1pt]{\scriptsize déploiements plus larges}};
        \node[anchor=west,text width=3.9cm,align=left] at (0.5,-3.6)  {{\bfseries Grand public}\\[-1pt]{\scriptsize familles, voyageurs, une fois le réseau visible}};
      \end{tikzpicture}
    \end{column}
    \begin{column}{0.25\textwidth}
      \begin{idee}\footnotesize\centering
        {\sffamily\bfseries\color{insarouge}L'objectif}\\[3pt]
        Créer d'abord une zone où le produit est vraiment utile…\\[3pt]
        …puis utiliser cette densité, les retours terrain et la preuve d'usage pour élargir.
      \end{idee}
    \end{column}
  \end{columns}
  \vspace{0.05cm}
  \begin{pcompact}\centering\footnotesize
    Ne pas viser «\,tout le monde\,» au lancement~: gagner un premier territoire d'usage, puis élargir.
  \end{pcompact}
  \srcnote{Choix du marché de départ d'après le SWOT de l'équipe~; les seuils de décision sont dans la partie Étude quantitative.}
\end{frame}

%  Fichier de contenu : ni préambule ni \begin{document}, commence par \section.
%
%  Remplace l'ancien SWOT en PDF (\includepdf supprimé) : le diagnostic est maintenant écrit
%  en LaTeX/TikZ, avec la charte du reste de la présentation, et complété par un SWOT croisé.
%  Ce qui a changé par rapport à la version PDF :
%   - la décomposition des 69 € n'est plus répétée ici (elle est dans la partie Finances) ;
%   - chaque point du diagnostic renvoie à un « frein » F1 à F7 de l'étude documentaire ;
%   - les marges des offres Solo / Duo / Groupe sont recalculées (marge par commande ET par kit,
%     avec le coût de vente de 4 € compté une seule fois par commande) ;
%   - le SWOT croisé (stratégies SO, WO, ST, WT) relie le diagnostic aux trois enjeux.
%  Hypothèses chiffrées : celles de la partie Finances (kit 69 € TTC, matériel 30,30 €, vente 4 €
%  par commande, 1 kit offert pour 10, charges fixes de l'année 3 = 430 k€).
%  Les offres Solo (79 €), Duo (139 €, 2 kits) et Groupe (249 €, 4 kits) sont celles du SWOT d'origine ;
%  le nombre de kits du Groupe (4) est une hypothèse.
%%%%%%%%%%%%%%%%%%%%%%%%%%%%%%%%%%%%%%%%%%%%%%%%%%%%%%%%%%%%%%%%%%%%%%%%%%%%%%
\section{SWOT et enjeux}

% Entrée de SWOT : \swit[couleur du code]{code}{titre}{texte}  (titre et texte sur le même paragraphe)
\newcommand{\swit}[4][insarouge]{%
  \begin{tcolorbox}[colback=insabeige,colframe=insabeige,boxrule=0pt,arc=2pt,left=4pt,right=4pt,
                    top=1.5pt,bottom=1.5pt,before skip=0pt,after skip=2pt]
    \scriptsize\raggedright{\sffamily\bfseries\color{#1}#2}\hspace{0.3em}{\sffamily\bfseries\footnotesize #3}\hspace{0.3em}#4\par
  \end{tcolorbox}}
% En-tête de colonne : \swhead{Forces}{couleur}
\newcommand{\swhead}[2]{%
  {\sffamily\bfseries\footnotesize\color{#2}\MakeUppercase{#1}}\\[-3pt]{\color{#2}\rule{\linewidth}{0.6pt}}\\[-1pt]}
% renvoi vers un enjeu ou un frein
\newcommand{\swref}[1]{\textcolor{insarouge}{\textbf{→\,#1}}}

%---------------------------------------------------------------------------
\begin{frame}{Diagnostic interne~: forces et faiblesses}
  \begin{columns}[T]
    \begin{column}{0.485\textwidth}
      \swhead{Forces}{insarouge}
      \swit{S1}{Proposition de valeur claire}{Une messagerie texte sans réseau mobile, sans compte ni abonnement.}
      \swit{S2}{Coût d'entrée faible}{69\,€ une fois, contre ≈ 830\,€ sur trois ans pour une balise satellite.}
      \swit{S3}{Écosystème libre existant}{Meshtastic, Heltec V3, applications~: l'effort porte sur l'usage.}
      \swit{S4}{Simplicité comme différence}{Réglages Europe, clé privée, QR code, guides en français.}
      \swit{S5}{Un socle, plusieurs usages}{Outdoor, campagne, ONG, crise~: un kit, quatre terrains.}
    \end{column}
    \begin{column}{0.485\textwidth}
      \swhead{Faiblesses}{insagris}
      \swit[insagris]{W1}{Faible barrière technologique}{Un bricoleur refait un kit proche pour 18 à 30\,€. \swref{F5}}
      \swit[insagris]{W2}{Effet de réseau indispensable}{Un kit isolé ne sert à rien~: la valeur vient des voisins. \swref{enjeu 1}}
      \swit[insagris]{W3}{Marge de contribution fine}{20,17\,€ par kit, mais 0,62\,€ de résultat par kit en année 3. \swref{enjeu 2}}
      \swit[insagris]{W4}{Financement avant rentabilité}{≈ 450\,k€ à financer, point mort ≈ 21\,300 kits par an. \swref{F7}}
      \swit[insagris]{W5}{Positionnement trop large}{ONG, randonneurs, makers, public~: attentes différentes. \swref{enjeu 3}}
    \end{column}
  \end{columns}
  \vspace{-0.2cm}
  \begin{pcompact}\centering\footnotesize
    Lecture clé~: la technologie est accessible~; la valeur viendra de l'usage, de la marque et de l'écosystème.
  \end{pcompact}
  \srcnote{Synthèse des parties Concurrence, Finances et Étude documentaire. F2 à F7 = freins de l'étude documentaire.}
\end{frame}

%---------------------------------------------------------------------------
\begin{frame}{Diagnostic externe~: opportunités et menaces}
  \begin{columns}[T]
    \begin{column}{0.485\textwidth}
      \swhead{Opportunités}{insarouge}
      \swit{O1}{Un marché outdoor important}{27 millions de pratiquants, dont 240\,000 licenciés FFRandonnée.}
      \swit{O2}{Besoin croissant de résilience}{Zones blanches, crises~: l'État invite à préparer un kit de 72\,h.}
      \swit{O3}{Réseaux locaux à construire}{Refuges, clubs, communes, ADRASEC~: points d'ancrage du maillage.}
      \swit{O4}{Précommandes adaptées}{Le financement participatif teste la demande et paie le premier stock.}
      \swit{O5}{Made in France, souveraineté}{Relocalisation par étapes~: un argument pour les collectivités.}
    \end{column}
    \begin{column}{0.485\textwidth}
      \swhead{Menaces}{insagris}
      \swit[insagris]{T1}{Fait-maison, kits concurrents}{Cartes dès 18\,€, kits prêts à l'emploi 59 à 153\,€. \swref{C}}
      \swit[insagris]{T2}{Substituts puissants}{SOS satellite des smartphones, balises, SMS satellite. \swref{F2}}
      \swit[insagris]{T3}{Contraintes réglementaires}{Marquage CE, cybersécurité, responsabilité produit. \swref{F3, F4}}
      \swit[insagris]{T4}{Dépendance amont}{Un seul fabricant de carte~; cap et marque de Meshtastic. \swref{F5, F6}}
      \swit[insagris]{T5}{Promesse mal comprise}{Ni téléphone, ni balise de détresse~: cadrer la promesse. \swref{F3}}
    \end{column}
  \end{columns}
  \vspace{-0.22cm}
  \begin{pcompact}\centering\footnotesize
    Opportunité réelle, mais fenêtre concurrentielle courte~: il faut bâtir vite une communauté et une marque.
  \end{pcompact}
  \srcnote{Sources~: \cite{ffrando,ffrando_licencies,guide_tous_resp}. F2 à F6 = freins de l'étude documentaire~; C = chapitre Concurrence.}
\end{frame}

%---------------------------------------------------------------------------
\newenvironment{pucesx}
  {\begingroup\setbeamerfont{itemize/enumerate body}{size=\scriptsize}\begin{itemize}\setlength{\itemsep}{0pt}\setlength{\topsep}{1pt}}
  {\end{itemize}\endgroup}
\begin{frame}{SWOT croisé~: quelles stratégies en déduire~?}
  \footnotesize\renewcommand{\arraystretch}{1.0}\setlength{\tabcolsep}{4pt}
  \hyphenpenalty=10000 \exhyphenpenalty=10000
  \noindent\begin{tabularx}{\textwidth}{@{}>{\raggedright\arraybackslash}p{1.65cm}>{\raggedright\arraybackslash}X>{\raggedright\arraybackslash}X@{}}
    & \cellcolor{insarose}\textcolor{insarouge}{\bfseries Opportunités (O)} & \cellcolor{insarose}\textcolor{insarouge}{\bfseries Menaces (T)} \\[2pt]
    \cellcolor{insarose}\textcolor{insarouge}{\bfseries Forces (S)} &
      \cellcolor{insabeige}{\sffamily\bfseries\color{insarouge}Attaquer}
      \begin{pucesx}
        \item Micro-réseaux prêts à l'emploi (Duo, Groupe) pour clubs, refuges, groupes de randonnée.
        \item Précommandes~: tester la demande et financer le premier stock.
        \item Pilotes avec des ADRASEC et des communes~: preuves d'usage, premiers nœuds.
      \end{pucesx} &
      \cellcolor{insabeige}{\sffamily\bfseries\color{insarouge}Défendre}
      \begin{pucesx}
        \item Vendre le prêt-à-l'emploi~: réglages, clé privée, guide, garantie, marquage CE.
        \item Rester un complément du smartphone et de la balise, sans promesse de secours.
        \item Contribuer au code libre, rester « compatible Meshtastic ».
      \end{pucesx} \\[3pt]
    \cellcolor{insarose}\textcolor{insarouge}{\bfseries Faiblesses (W)} &
      \cellcolor{insabeige}{\sffamily\bfseries\color{insarouge}Corriger}
      \begin{pucesx}
        \item Densifier autour de points d'ancrage (refuge, club, mairie) contre l'effet de réseau.
        \item Tester Solo, Duo et Groupe avant la série pour relever la marge.
        \item Bâtir le dossier de financement sur les précommandes.
      \end{pucesx} &
      \cellcolor{insabeige}{\sffamily\bfseries\color{insarouge}Éviter}
      \begin{pucesx}
        \item Notice claire, test à 100\,\%, assurance responsabilité produit.
        \item Double source de cartes (V3 et V4), stock tampon, vente d'abord en France.
        \item Fixer des seuils d'arrêt avant le moule et la série (étude quantitative).
      \end{pucesx} \\
  \end{tabularx}
  \srcnote{Croisement des deux diapositives précédentes (méthode TOWS)~: S, W = forces et faiblesses~; O, T = opportunités et menaces.}
\end{frame}

%---------------------------------------------------------------------------
% Liaisons possibles entre n kits : n(n-1)/2 (2 kits = 1, 4 = 6, 6 = 15, 10 = 45)
\begin{frame}{Enjeu 1~: résoudre l'effet de réseau}
  \begin{columns}[T]
    \begin{column}{0.46\textwidth}
      \centering
      \begin{tikzpicture}[font=\footnotesize]
        \node[pcli,text width=2.3cm,minimum height=0.7cm] (a) at (0,1.05) {Peu d'utilisateurs};
        \node[pgris,text width=1.6cm,minimum height=0.7cm] (b) at (3.55,1.05) {Peu de relais};
        \node[pgris,text width=1.6cm,minimum height=0.7cm] (c) at (1.78,-0.3) {Valeur limitée};
        \draw[fluxr] (a) -- (b);
        \draw[fluxr] (b) -- (c);
        \draw[fluxr] (c) -- (a);
        \node[font=\scriptsize\bfseries,text=insarouge,align=center] at (1.78,0.5) {cercle\\vicieux};
      \end{tikzpicture}

      \vspace{0.05cm}
      {\footnotesize\textbf{Un groupe vaut plus que la somme de ses kits}~: $n$ kits offrent $n(n{-}1)/2$ liaisons.\par}
      \begin{tikzpicture}[font=\scriptsize]
        \tikzset{noeud/.style={circle,fill=insarouge,inner sep=1.3pt}}
        % 2 kits
        \begin{scope}[shift={(0.0,0)}]
          \node[noeud] (u0) at (-0.28,0) {};\node[noeud] (u1) at (0.28,0) {};
          \draw[insagris,thick] (u0) -- (u1);
          \node[align=center,anchor=north] at (0,-0.55) {2 kits\\\textbf{1 liaison}};
        \end{scope}
        % 4 kits
        \begin{scope}[shift={(1.85,0)}]
          \foreach \i in {0,...,3} {\node[noeud] (v\i) at ({45+90*\i}:0.4) {};}
          \foreach \i/\j in {0/1,0/2,0/3,1/2,1/3,2/3} {\draw[insagris,thick] (v\i) -- (v\j);}
          \node[align=center,anchor=north] at (0,-0.55) {4 kits\\\textbf{6 liaisons}};
        \end{scope}
        % 6 kits
        \begin{scope}[shift={(3.7,0)}]
          \foreach \i in {0,...,5} {\node[noeud] (w\i) at ({90+60*\i}:0.42) {};}
          \foreach \i/\j in {0/1,0/2,0/3,0/4,0/5,1/2,1/3,1/4,1/5,2/3,2/4,2/5,3/4,3/5,4/5} {\draw[insagris,thin] (w\i) -- (w\j);}
          \node[align=center,anchor=north] at (0,-0.55) {6 kits\\\textbf{15 liaisons}};
        \end{scope}
        \node[align=center,font=\scriptsize] at (5.15,0) {10 kits\\\textbf{45}};
      \end{tikzpicture}
    \end{column}
    \begin{column}{0.51\textwidth}
      \begin{columns}[T]
        \begin{column}{0.47\linewidth}
          \begin{carte}[height=1.7cm,valign=top]
            \centering
            \begin{tikzpicture}\path (-0.5,-0.35) rectangle (0.5,0.35);\pic[insarouge,scale=0.9] at (0,0) {icone personne};\end{tikzpicture}\\[1pt]
            \gros{Duo, famille}\\[1pt]{\footnotesize 2 à 3 appareils}
          \end{carte}
        \end{column}
        \begin{column}{0.47\linewidth}
          \begin{carte}[height=1.7cm,valign=top]
            \centering
            \begin{tikzpicture}\path (-0.5,-0.35) rectangle (0.5,0.35);\pic[scale=0.9] at (0,0) {icone mont};\end{tikzpicture}\\[1pt]
            \gros{Randonnée}\\[1pt]{\footnotesize 4 appareils}
          \end{carte}
        \end{column}
      \end{columns}
      \vspace{0.1cm}
      \begin{columns}[T]
        \begin{column}{0.47\linewidth}
          \begin{carte}[height=1.7cm,valign=top]
            \centering
            \begin{tikzpicture}\path (-0.5,-0.35) rectangle (0.5,0.35);\pic[insarouge,scale=0.9] at (0,0) {icone tente};\end{tikzpicture}\\[1pt]
            \gros{Refuge, club}\\[1pt]{\footnotesize un relais et ses utilisateurs}
          \end{carte}
        \end{column}
        \begin{column}{0.47\linewidth}
          \begin{carte}[height=1.7cm,valign=top]
            \centering
            \begin{tikzpicture}\path (-0.5,-0.35) rectangle (0.5,0.35);\pic[scale=0.9] at (0,0) {icone croix};\end{tikzpicture}\\[1pt]
            \gros{ONG}\\[1pt]{\footnotesize secours~: 10 à 30 appareils}
          \end{carte}
        \end{column}
      \end{columns}
    \end{column}
  \end{columns}
  \vspace{0.05cm}
  \begin{pcompact}\centering\footnotesize
    Le bon produit est peut-être moins «\,une antenne\,» qu'un réseau prêt à l'emploi pour un groupe constitué.
  \end{pcompact}
  \srcnote{Micro-réseaux d'après le SWOT de l'équipe~; tailles de groupe à mesurer dans le questionnaire (questions 6 et 18).}
\end{frame}

%---------------------------------------------------------------------------
\begin{frame}{Enjeu 2~: un prix d'appel, une marge à protéger}
  \footnotesize\renewcommand{\arraystretch}{1.3}
  \hyphenpenalty=10000 \exhyphenpenalty=10000
  \newcolumntype{N}{>{\raggedleft\arraybackslash}p{1.5cm}}
  \setlength{\tabcolsep}{3pt}
  \noindent\begin{tabular}{@{}>{\raggedright\arraybackslash}p{2.9cm}NNNNNN@{}}
    \rowcolor{insarose}\textcolor{insarouge}{\bfseries Offre} & \textcolor{insarouge}{\bfseries Prix TTC} & \textcolor{insarouge}{\bfseries Par kit} & \textcolor{insarouge}{\bfseries Marge par commande} & \textcolor{insarouge}{\bfseries Marge par kit} & \textcolor{insarouge}{\bfseries Point mort (kits/an)} & \textcolor{insarouge}{\bfseries Résultat année 3}\\
    Kit seul (référence)    & 69\,€  & 69,00\,€ & 20,17\,€ & 20,17\,€ & 21\,300 & +13,7\,k€\\
    Solo                    & 79\,€  & 79,00\,€ & 28,50\,€ & 28,50\,€ & 15\,100 & +197\,k€\\
    \rowcolor{insabeige}Duo (2 kits)           & 139\,€ & 69,50\,€ & 45,17\,€ & 22,59\,€ & 19\,000 & +67\,k€\\
    Groupe (4 kits)         & 249\,€ & 62,25\,€ & 70,18\,€ & 17,54\,€ & 24\,500 & −44\,k€\\
  \end{tabular}
  \vspace{0.12cm}
  \begin{columns}[T]
    \begin{column}{0.325\textwidth}
      \begin{carte}[height=1.7cm,valign=top]
        \raggedright\gros{−73\,k€}\\[1pt]{\footnotesize sur le résultat de l'année 3 si le coût matériel monte de 10\,\% (+3\,€ par kit).}
      \end{carte}
    \end{column}
    \begin{column}{0.325\textwidth}
      \begin{carte}[height=1.7cm,valign=top]
        \raggedright\gros{−44\,k€}\\[1pt]{\footnotesize si la vente, le SAV et la logistique coûtent 2\,€ de plus par commande.}
      \end{carte}
    \end{column}
    \begin{column}{0.325\textwidth}
      \begin{carte}[height=1.7cm,valign=top]
        \raggedright\gros{±92\,k€}\\[1pt]{\footnotesize pour 5\,€ de prix en plus ou en moins sur chaque kit.}
      \end{carte}
    \end{column}
  \end{columns}
  \vspace{0.08cm}
  \begin{pcompact}\centering\footnotesize
    Le Duo relève la marge par kit~; le Groupe à 249\,€ la baisse (il faudrait ≥ 262\,€). Le prix se tranche par le test (questions 19 à 21 du questionnaire), pas par intuition.
  \end{pcompact}
  \srcnote{Calcul avec le modèle de la partie Finances~: matériel 30,30\,€, vente 4\,€ par commande, 1 kit offert pour 10, charges fixes 430\,k€, 22\,000 kits. Marge = prix HT moins coûts variables.}
\end{frame}

%---------------------------------------------------------------------------
\begin{frame}{Enjeu 3~: choisir un marché de départ}
  \begin{columns}[T]
    \begin{column}{0.35\textwidth}
      \begin{carte}[height=4.45cm,valign=top]
        \raggedright
        {\sffamily\bfseries\scriptsize\color{insarouge}CIBLE DE DÉPART (HYPOTHÈSE)}\\[2pt]
        {\sffamily\bfseries\large\color{insarouge}Randonneurs et groupes outdoor}\\[2pt]
        {\footnotesize\textbf{Pourquoi~?}}
        \begin{puces}
          \item besoin simple à comprendre
          \item usage réel en zone blanche
          \item marché large et accessible
          \item vente plus simple que le secours
        \end{puces}
        {\scriptsize\textcolor{insarouge}{À confirmer par l'étude quantitative.}}
      \end{carte}
    \end{column}
    \begin{column}{0.36\textwidth}
      {\sffamily\bfseries\footnotesize Progression naturelle}\\[3pt]
      \begin{tikzpicture}[font=\footnotesize]
        \draw[insarouge!40,thick] (0.2,0.05) -- (0.2,-3.55);
        \node[num] at (0.2,0)     {1};
        \node[num] at (0.2,-0.9) {2};
        \node[num] at (0.2,-1.8)  {3};
        \node[num] at (0.2,-2.7) {4};
        \node[num] at (0.2,-3.6)  {5};
        \node[anchor=west,text width=3.9cm,align=left] at (0.5,0)     {{\bfseries Outdoor}\\[-1pt]{\scriptsize randonneurs et groupes}};
        \node[anchor=west,text width=3.9cm,align=left] at (0.5,-0.9) {{\bfseries Refuges et clubs}\\[-1pt]{\scriptsize points de relais locaux}};
        \node[anchor=west,text width=3.9cm,align=left] at (0.5,-1.8)  {{\bfseries Associations}\\[-1pt]{\scriptsize ADRASEC, communautés structurées}};
        \node[anchor=west,text width=3.9cm,align=left] at (0.5,-2.7) {{\bfseries ONG, communes, pros}\\[-1pt]{\scriptsize déploiements plus larges}};
        \node[anchor=west,text width=3.9cm,align=left] at (0.5,-3.6)  {{\bfseries Grand public}\\[-1pt]{\scriptsize familles, voyageurs, une fois le réseau visible}};
      \end{tikzpicture}
    \end{column}
    \begin{column}{0.25\textwidth}
      \begin{idee}\footnotesize\centering
        {\sffamily\bfseries\color{insarouge}L'objectif}\\[3pt]
        Créer d'abord une zone où le produit est vraiment utile…\\[3pt]
        …puis utiliser cette densité, les retours terrain et la preuve d'usage pour élargir.
      \end{idee}
    \end{column}
  \end{columns}
  \vspace{0.05cm}
  \begin{pcompact}\centering\footnotesize
    Ne pas viser «\,tout le monde\,» au lancement~: gagner un premier territoire d'usage, puis élargir.
  \end{pcompact}
  \srcnote{Choix du marché de départ d'après le SWOT de l'équipe~; les seuils de décision sont dans la partie Étude quantitative.}
\end{frame}
